%%% Hlavní soubor. Zde se definují základní parametry a odkazuje se na ostatní části. %%%

%% Verze pro jednostranný tisk:
% Okraje: levý 40mm, pravý 25mm, horní a dolní 25mm
% (ale pozor, LaTeX si sám přidává 1in)
\documentclass[hidelinks,12pt,a4paper]{report}
\setlength\textwidth{145mm}
\setlength\textheight{247mm}
\setlength\oddsidemargin{15mm}
\setlength\evensidemargin{15mm}
\setlength\topmargin{0mm}
\setlength\headsep{0mm}
\setlength\headheight{0mm}
% \openright zařídí, aby následující text začínal na pravé straně knihy
\let\openright=\clearpage

%% Pokud tiskneme oboustranně:
% \documentclass[12pt,a4paper,twoside,openright]{report}
% \setlength\textwidth{145mm}
% \setlength\textheight{247mm}
% \setlength\oddsidemargin{14.2mm}
% \setlength\evensidemargin{0mm}
% \setlength\topmargin{0mm}
% \setlength\headsep{0mm}
% \setlength\headheight{0mm}
% \let\openright=\cleardoublepage

%% Vytváříme PDF/A-2u
\usepackage[a-2u]{pdfx}

%% Přepneme na českou sazbu a fonty Latin Modern
\usepackage[czech]{babel}
\usepackage{lmodern}
\usepackage[T1]{fontenc}
\usepackage{textcomp}

%% Použité kódování znaků: obvykle latin2, cp1250 nebo utf8:
\usepackage[utf8]{inputenc}

%%% Další užitečné balíčky (jsou součástí běžných distribucí LaTeXu)
\usepackage{amsmath}        % rozšíření pro sazbu matematiky
\usepackage{amsfonts}       % matematické fonty
\usepackage{amsthm}         % sazba vět, definic apod.
\usepackage{bbding}         % balíček s nejrůznějšími symboly
			    % (čtverečky, hvězdičky, tužtičky, nůžtičky, ...)
\usepackage{bm}             % tučné symboly (příkaz \bm)
\usepackage{graphicx}       % vkládání obrázků
\usepackage{fancyvrb}       % vylepšené prostředí pro strojové písmo
\usepackage{indentfirst}    % zavede odsazení 1. odstavce kapitoly
\usepackage{natbib}         % zajištuje možnost odkazovat na literaturu
			    % stylem AUTOR (ROK), resp. AUTOR [ČÍSLO]
\usepackage[nottoc]{tocbibind} % zajistí přidání seznamu literatury,
                            % obrázků a tabulek do obsahu
\usepackage{icomma}         % inteligetní čárka v matematickém módu
\usepackage{dcolumn}        % lepší zarovnání sloupců v tabulkách
\usepackage{booktabs}       % lepší vodorovné linky v tabulkách
\usepackage{paralist}       % lepší enumerate a itemize
\usepackage{xcolor}         % barevná sazba
\usepackage{multicol}

\usepackage{rotating}
%\usepackage{fontspec}
\usepackage{tipa}
\usepackage{enumitem}
\usepackage{makecell}
\usepackage{alltt}
\usepackage{refcount}
\usepackage{tabularx}
\usepackage{caption}

%\usepackage{etoolbox}
%\makeatletter
%\patchcmd{\chapter}{\if@openright\cleardoublepage\else\clearpage\fi}{}{}{}
%\makeatother


%%% Údaje o práci

% Název práce v jazyce práce (přesně podle zadání)
\def\NazevPrace{Reflexe vybraných témat křesťanského myšlení v pojetí Karla Makoně}

% Název práce v angličtině
\def\NazevPraceEN{Reflecting Selected Topics in Christian Thinking in Karel Makoň's Conception}

% Jméno autora
\def\AutorPrace{RNDr. Jan Oldřich Krůza, Ph.D.}

% Rok odevzdání
\def\RokOdevzdani{2022}

% Název katedry nebo ústavu, kde byla práce oficiálně zadána
% (dle Organizační struktury MFF UK, případně plný název pracoviště mimo MFF)
\def\Katedra{Katedra systematické teologie, teologické etiky a teologické filozofie}
\def\KatedraEN{Department of Systematical Theology, Theological Ethics and Theological Philosophy}

% Jedná se o katedru (department) nebo o ústav (institute)?
\def\TypPracoviste{Katedra}
\def\TypPracovisteEN{Department}

% Vedoucí práce: Jméno a příjmení s~tituly
\def\Vedouci{doc. ThDr. Jiří Vogel, Th.D.}

% Pracoviště vedoucího (opět dle Organizační struktury MFF)
\def\KatedraVedouciho{\Katedra}
\def\KatedraVedoucihoEN{\KatedraEN}

% Studijní program a obor
\def\StudijniProgram{husitská teologie}
\def\StudijniObor{HT NK}

% Nepovinné poděkování (vedoucímu práce, konzultantovi, tomu, kdo
% zapůjčil software, literaturu apod.)
%\def\Podekovani{%
%Děkuji
%}

% Abstrakt (doporučený rozsah cca 80-200 slov; nejedná se o zadání práce)
\def\Abstrakt{%
Český mystik Karel Makoň (*1912 \textdagger1993) po sobě zanechal rozsáhlé a významné dílo v~psané a
mluvené podobě.
%Po mystickém zážitku v~koncentračním táboře strávil zbytek života vedením ostatních k~zasvěcení života hledání Království Božího.
V~desítkách
knih, které napsal, rozvádí návod k~využití pozemského života pro vstup do
vědomého spojení s~Věčností. Hovořil ke skupinkám příznivců, z~čehož se
zachovalo přes tisíc hodin magnetofonových záznamů.
Pro šíři jeho hovorů je těžké určit, jaké jsou hlavní body jeho poselství, a
dostupná literatura o něm je velmi sporadická.
Vyskytují se závažné výroky, které v~obměnách opakuje po více
než dvacet let svého auditivně zaznamenaného působení. Mezi tyto opakující se
výroky patří například: ,,Trpnost je nezbytnou součástí cesty k~Bohu.``,
,,Modlitba nesmí být mechanická, aby měla spojovací účinek.``, ,,Milost Boží je
zákonitým jevem.``
Tyto lze s~jistotou zařadit mezi stěžejní pilíře jeho nauky.
Jaké jsou ale ty ostatní?
Pomocí počítačové analýzy přepisů nahrávek lze dospět k~seznamu kandidátských témat,
jenž může být základem pro vyčerpávající seznam pokrývající celý mluvený korpus.
V~mluveném korpusu Karla Makoně lze nalézt odpovědi na všechny základní otázky systematické teologie.
Mnohé z~těchto odpovědí jsou v~kontroverzi k~církevním naukám.
Esence Makoňova poselství lze shrnout starověkým citátem ,,Tento život je mostem do věčnosti.``
Jiní jeho příznivci uvádějí ale jiná shrnutí.
}

\def\AbstraktEN{%
Karel Makoň (1912-1993) has left behind an extensive and significant opus, both
written and spoken. After a mystical experience in a concentration camp, he
spent the rest of his life urging others to devote their lives to seeking the
Kingdom of God. In the dozens of books he has written he elaborates on a how-to
for entering the conscious Life Eternal. He was talking to several groups of
followers, which resulted in over a thousand hours of magnetophone tape
recordings. Available literature about him is very sporadical. The breadth of
his talks makes it hard to determine the main points of his message. There are
repeated significant statements appearing over the span of more than twenty
years of his recorded oral activity. Some of these repeated statements are:
``Passivity is an essential part of the path to God.'', ``Prayer must not be
mechanical if it is to have its connecting effect.'', ``God's grace is an
inevitable phenomenon.'' A list of candidate topics in Makoň's spoken corpus can
be obtained through computational analysis of his transcribed recordings. This
can be a basis for an exhaustive list of topics covering the entire material.
All foundational questoins of systematical theology can be answered using
Makoň's talks. Many of these answers are controversial towards ecclesiastical
doctrines. The essence of Makoň's message can be summed up using the ancient
quote ``The life is a bridge into the Eternity''. Others come up with different
summaries though.
}


% 3 až 5 klíčových slov (doporučeno), každé uzavřeno ve složených závorkách
\def\KlicovaSlova{%
{Karel Makoň}, {Karel Weinfurter}, {křesťanská mystika}, {křesťanství a indická filozofie}
}
\def\KlicovaSlovaEN{%
{Karel Makoň}, {Karel Weinfurter}, {Christian Mysticism}, {Chstianity and Indian Philosophy}
}

%% Balíček hyperref, kterým jdou vyrábět klikací odkazy v PDF,
%% ale hlavně ho používáme k uložení metadat do PDF (včetně obsahu).
%% Většinu nastavítek přednastaví balíček pdfx.
\hypersetup{unicode}
\hypersetup{breaklinks=true}

%% Definice různých užitečných maker (viz popis uvnitř souboru)
%%% Tento soubor obsahuje definice různých užitečných maker a prostředí %%%
%%% Další makra připisujte sem, ať nepřekáží v ostatních souborech.     %%%

%%% Drobné úpravy stylu

% Tato makra přesvědčují mírně ošklivým trikem LaTeX, aby hlavičky kapitol
% sázel příčetněji a nevynechával nad nimi spoustu místa. Směle ignorujte.
\makeatletter
\def\@makechapterhead#1{
  {\parindent \z@ \raggedright \normalfont
   \Huge\bfseries \thechapter. #1
   \par\nobreak
   \vskip 20\p@
}}
\def\@makeschapterhead#1{
  {\parindent \z@ \raggedright \normalfont
   \Huge\bfseries #1
   \par\nobreak
   \vskip 20\p@
}}
\makeatother

% Toto makro definuje kapitolu, která není očíslovaná, ale je uvedena v obsahu.
\def\chapwithtoc#1{
\chapter*{#1}
\addcontentsline{toc}{chapter}{#1}
}

% Trochu volnější nastavení dělení slov, než je default.
\lefthyphenmin=2
\righthyphenmin=2

% Zapne černé "slimáky" na koncích řádků, které přetekly, abychom si
% jich lépe všimli.
%\overfullrule=1mm

%%% Makra pro definice, věty, tvrzení, příklady, ... (vyžaduje baliček amsthm)

\theoremstyle{plain}
\newtheorem{veta}{Věta}
\newtheorem{lemma}[veta]{Lemma}
\newtheorem{tvrz}[veta]{Tvrzení}

\theoremstyle{plain}
\newtheorem{definice}{Definice}

\theoremstyle{remark}
\newtheorem*{dusl}{Důsledek}
\newtheorem*{pozn}{Poznámka}
\newtheorem*{prikl}{Příklad}

%%% Prostředí pro důkazy

\newenvironment{dukaz}{
  \par\medskip\noindent
  \textit{Důkaz}.
}{
\newline
\rightline{$\qedsymbol$}
}

%%% Prostředí pro sazbu kódu, případně vstupu/výstupu počítačových
%%% programů. (Vyžaduje balíček fancyvrb -- fancy verbatim.)

\DefineVerbatimEnvironment{code}{Verbatim}{fontsize=\small, frame=single}

%%% Prostor reálných, resp. přirozených čísel
\newcommand{\R}{\mathbb{R}}
\newcommand{\N}{\mathbb{N}}

%%% Užitečné operátory pro statistiku a pravděpodobnost
\DeclareMathOperator{\pr}{\textsf{P}}
\DeclareMathOperator{\E}{\textsf{E}\,}
\DeclareMathOperator{\var}{\textrm{var}}
\DeclareMathOperator{\sd}{\textrm{sd}}

%%% Příkaz pro transpozici vektoru/matice
\newcommand{\T}[1]{#1^\top}

%%% Vychytávky pro matematiku
\newcommand{\goto}{\rightarrow}
\newcommand{\gotop}{\stackrel{P}{\longrightarrow}}
\newcommand{\maon}[1]{o(n^{#1})}
\newcommand{\abs}[1]{\left|{#1}\right|}
\newcommand{\dint}{\int_0^\tau\!\!\int_0^\tau}
\newcommand{\isqr}[1]{\frac{1}{\sqrt{#1}}}

%%% Vychytávky pro tabulky
\newcommand{\pulrad}[1]{\raisebox{1.5ex}[0pt]{#1}}
\newcommand{\mc}[1]{\multicolumn{1}{c}{#1}}

\newcommand{\ridv}[3]{\href{http://radio.makon.cz/zaznam/#1\#ts=#3}{\texttt{#1} od pozice #2}}
\newcommand{\rid}[3]{\href{http://radio.makon.cz/zaznam/#1\#ts=#3}{\texttt{#1} #2}}
\newcommand{\rref}[3]{\footnote{\rid{#1}{#2}{#3}}}


%% Titulní strana a různé povinné informační strany
\begin{document}
%%% Titulní strana práce a další povinné informační strany

%%% Titulní strana práce

\pagestyle{empty}
\hypersetup{pageanchor=false}

\begin{center}

{\bf\Large UNIVERZITA KARLOVA}

\vfill

{\Large HUSITSKÁ TEOLOGICKÁ FAKULTA}

\vfill

\centerline{\mbox{\includegraphics[width=66mm]{rc/logo-htf.jpg}}}

\vfill

{\bfseries\NazevPrace}

\vfill

{\bfseries\NazevPraceEN}

\vfill

Diplomová práce

\vspace{15mm}

\begin{multicols}{2}
Vedoucí práce:\\
\Vedouci\\
\columnbreak
Autor:\\
\AutorPrace\\
\end{multicols}

\vfill

% Zde doplňte rok
Praha \RokOdevzdani

\end{center}

\newpage

%%% Strana s čestným prohlášením k disertační práci

\openright
\hypersetup{pageanchor=true}
\pagestyle{plain}
\pagenumbering{roman}
\vglue 0pt plus 1fill

\noindent
Prohlašuji, že jsem tuto seminární práci vypracoval samostatně a výhradně
s~použitím citovaných pramenů, literatury a dalších odborných zdrojů.

\medskip\noindent
Beru na~vědomí, že se na moji práci vztahují práva a povinnosti vyplývající
ze zákona č. 121/2000 Sb., autorského zákona v~platném znění, zejména skutečnost,
že Univerzita Karlova má právo na~uzavření licenční smlouvy o~užití této
práce jako školního díla podle §60 odst. 1 autorského zákona.

\vspace{10mm}

\hbox{\hbox to 0.5\hsize{%
%V \hbox to 6em{\dotfill} dne \hbox to 6em{\dotfill}
V Praze dne \today
\hss}\hbox to 0.5\hsize{\dotfill\quad}}
\smallskip
\hbox{\hbox to 0.5\hsize{}\hbox to 0.5\hsize{\hfil Podpis autora\hfil}}

\vspace{20mm}
\newpage

%%% Poděkování

\openright

\noindent
%\Podekovani

\newpage

%%% Povinná informační strana disertační práce

\openright

\vbox to 0.5\vsize{
\setlength\parindent{0mm}
\setlength\parskip{5mm}

{\bf Abstrakt:\\}
\Abstrakt

{\bf Klíčová slova:\\}
\KlicovaSlova

\vfill

{\bf Abstract:\\}
\AbstraktEN

{\bf Keywords:\\}
\KlicovaSlovaEN

\vss}

\newpage



\openright


%%% Strana s automaticky generovaným obsahem disertační práce

\tableofcontents

\newpage
\pagenumbering{arabic}
\setcounter{page}{1}
%%% Jednotlivé kapitoly práce jsou pro přehlednost uloženy v samostatných souborech

\chapter{Úvod}
\label{kap:uvod}

Karel Makoň je téměř neznámou postavou české literatury a kupodivu i české
mystiky. Jeho knihy nejsou prakticky v~žádné knihovně o knihkupectvích nemluvě.
Není o něm zmínka v~encyklopediích, v~době psaní tohoto textu (duben 2022) ani
na Wikipedii. Jeho dílo je přitom jedinečné jak rozsahem (27 autorských spisů,
28 překladů a komentářů k~nim, přes 1000 hodin záznamů přednášek) tak i
závažností obsahu.

Když jsem se s~jeho dílem poprvé setkal v~roce 2008, uchvátilo mě natolik, že
jsem jeho zpracování bez nadsázky zasvětil život. Dva roky jsem digitalizoval
magnetofonové pásky se záznamy jeho promluv, pak jsem jedenáct let pracoval
v~rámci doktorského studia na MFF UK na jejich přepisu a zpřístupnění. Zpracovat
jeho dílo nejen z~hlediska informatického jakožto sadu dat, nýbrž z~hlediska
obsahu, bylo mým snem od momentu, kdy jsem si uvědomil, že v~rámci disertace
v~oboru komputační lingvistiky se musím více či méně držet mantinelů technického
zaměření pracoviště.

Rozsah Makoňova díla je takový, že celé je přečíst a poslechnout je na mnoho
let. Sám jsem se k~němu dostal nejprve přes nahrávky. Tato preference mi
zůstala. Poslechl jsem zatím přibližně třetinu Makoňových hovorů, zatímco knih
jsem přečetl jenom hrstku. Svoji disertační práci jsem též věnoval nahrávkám. Zatímco
knihy byly úsilím předchozích Makoňových příznivců vydávány, nahrávkám se
dosud nikdo systematicky nevěnoval. Proto tomuto aspektu Makoňovy tvorby zůstávám věren i
zde v~této diplomové práci.

Za její úkol jsem si vytyčil zodpovědět otázku \textit{,,Co vlastně říká Karel
Makoň?{}``} Soustředím se tedy především na záznamy Makoňových hovorů a knihy
zmiňuji jen okrajově. Samozřejmě by se dalo postupovat i zcela opačně a postavit
práci na knihách. Věřím tomu, že takové práce se v~dohlednu někdo ujme s~nemenší
vervou, než jsem se já ujal zpracování nahrávek. Kéž k~tomu tato práce pomůže.

\section{Citace}

Jelikož se tato práce zabývá Makoňovými nahrávkami, často se odkazuji na
konkrétní pasáže. Nahrávky nejsou běžným materiálem k~citování a těžko na ně lze
uplatnit standardní citační praxi. Soubor digitalizovaných nahrávek je
%coby výsledek mojí předchozí snahy
přístupný ze dvou hlavních kanálů:
\begin{enumerate}
\item{
  Z~internetové stránky \url{radio.makon.cz}, kde se nachází mnou vytvořená
  webová služba, která umožňuje vyhledávání a poslech se simultánním zobrazením
  přepisů a
}
\item{
  z~repozitáře Lindat/Clariah, který je výsledkem společného snažení několika
  institucí včetně Univerzity Karlovy a jejímž záměrem je otevřený přístup
  k~jazykovým zdrojům: \url{https://lindat.cz/repository/xmlui/handle/11234/1-3422}.
}
\end{enumerate}
Přístup z~webové aplikace je pohodlnější, ovšem u Lindat/Clariah je garantována
trvalost přístupu.

Citace z~Makoňova mluveného korpusu udávám ve formě párů \textit{identifikátor
nahrávky} - \textit{časová pozice}. Identifikátory vznikaly při digitalizaci tím
způsobem, že se opisovaly popisky na nosičích (hlavně kazetách a kotoučích). U
identifikátorů je garantováno, že neobsahují mezery, ale jinak jejich formát
nemá omezení. V~textu indentifikátor vždy uvádím pomocí monospacového písma, např.
\texttt{85-05A}. Pozice je od začátku nahrávky a to v~běžném formátu
,,\textit{minuty}:\textit{sekundy}``, popř
,,\textit{hodiny}:\textit{minuty}:\textit{sekundy}``.

\include{predstaveni-km}
\chapter{Závěr}

\section{Esence Makoňova poselství}

Nemělo by být překvapením, že po napsání této práce docházím ke shrnutí Makoňova
poselství v~tomto smyslu:
\begin{quote}
,,Smyslem lidského života je vědomé spojení s~Bohem,
což se rovná vědomému prožívání vlastní věčné podstaty. Vše má sloužit jako
prostředek k~tomuto cíli.``
\end{quote}
Elegantněji by věc šla vyjádřit slovy starověkého citátu:
\begin{quote}
,,Tento život je mostem do věčnosti.``
\end{quote}
Jakkoliv se může zpochybňovat jeho historická vazba na rané
křesťanství,\cite{dus2001neznama} jeho pravdivost byla pro Makoně natolik
strhující, že se stal jádrem jeho nauky.

Tento závěr je každopádně subjektivní, jakkoliv je podložen četnými ukázkami.
Abych si ověřil validitu svého závěru, poprosil jsem další příznivce Makoňovy
nauky, aby mi sdělili v~jedné či několika málo větách, co oni považují za esenci
Makoňova poselství. Každý z~dotázaných na rozdíl ode mne Makoně osobně zažil a
byli jím formováni po mnoho let svých životů. Zde jsou jejich odpovědi:

\begin{itemize}
  \item{\textit{,,Cíl vždycky nejvyšší.``}}
  \item{\textit{,,Nelpěte na událostech, aby dopadaly podle vašich představ. Nechte věci, ať
    se dějí.``}}
\item{\textit{,,Zapomínat na sebe. Jsme povinováni láskou ke všemu a ke všem. Nenávist
  není přípustná. I ten, kdo nejde cestou lásky, musí zapomínat sám na sebe, aby
    vystoupil ze svého já.``}}
  \item{\textit{,,Nepřekážejte Bohu.``}}
  \item{\textit{,,V~každém živém tvoru musí být Bůh neustále přítomen: ,Bohové jste.`
    Největší překážkou na cestě k~Bohu je vlastnění všeho druhu.``}}
\end{itemize}

Je vidět, že rozmanitost jednotlivých shrnutí vysoko převyšuje jejich
příbuznost. Je to zajímavá ukázka toho, jak je každý prakticky tímtéž ovlivněn
rozdílně dle svého založení. Každopádně lze s~radostí konstatovat, že jednotlivé
varianty si snad neodporují.

\section{Karel Makoň jako náboženská autorita}

Z~církevního hlediska se Karel Makoň může jevit jako kacíř -- a je-li kacířství odchýlení se od
církevní
nauky,\footnote{\href{https://www.oxfordreference.com/view/10.1093/oi/authority.20110803095932433}{oxfordreference.com/view/10.1093/oi/authority.20110803095932433}}
pak je kacířem zcela přiznaně a otevřeně. Sám tvrdí, že definicí
kacířství je vydávání části učení za celek, proto sebe za kacíře nepovažuje. Je
důležité vzít v~potaz, že Makoň se vůbec nesnaží zaštiťovat autoritou nějaké
církve. Nemluví jako křesťan, ale jako člověk, který čerpá z~vlastního vhledu.
Církev vidí často zkresleně, nebo aspoň dnešní církve neodpovídají tomu, jak o
nich Makoň mluví, což staví jejich vývoj za poslední dekády do velmi dobrého
světla. Makoň vidí v~katolické tradici nedoceněný a nesmírně propracovaný
prostředek na cestě k~vědomému spojení s~Bohem. Ostatně všechno na světě, i
Ježíše samotného a celý lidský život považuje za prostředek na cestě ke spojení
s~Bohem. Biblický text a katolickou tradici ale hodnotí jako obzvláštní
prostředky, protože nejsou tak kontextově závislé jako ostatní. Může se jimi
řídit každý a v~každé fázi.\rref{76-01-Kaly}{01:20:53} Biblickému textu je potřeba správně rozumět a
katolickou tradici je třeba náležitě opravit.

Z~katolického hlediska je tedy kacířem zcela jasně. Z~hlediska reformovaného
možná také. Ovšem zda má Makoň pravdu, to je zcela jiná otázka. Každopádně
slibuje něco, co církev neslibuje. Nejde mu o spásu po smrti, ani o víru v~to,
že už jsme spaseni. Jde mu o spásu reálnou, zcela jasně rozlišenou, u které
absolutně není pochyb, zda nastane či nastala. A celý život zasvěcuje tomu, aby
jí došli pokud možno všichni lidé. Věří, že to lze. Sám ji zažívá, spatřuje v~ní
jedinou skutečnou hodnotu a navádí k~ní s~nadlidskou vytrvalostí, přesvědčivostí
a mnohdy i zázračně.\footnote{Viz např. Mozaika mystického života Karla Makoně
str. 29 - 48.}

Mohl by se nám jevit jako egocentrik. Mluví sám o sobě a svých zážitcích častěji
než kdo jiný. Svoje zážitky bere jako model pro ostatní. Na některých místech je
patrno, že si je vědom svojí převahy. Vidět ho ale jako egoistu by byl
fatální omyl. Svoji výjimečnost zapřahá nekompromisně do svého poslání přimět
svoje posluchače, aby se stali dokonalými následovníky Kristovými. Říká, že
Satan je jáčko a radí ho přemoct tím, že ho zapojíme do spasitelského úkolu. Ať
má sám jáčko velké jakkoliv (a po pravdě, vzhledem k~jeho objektivním kvalitám,
se mně osobně zdá příkladně skromným), aplikuje na sobě tuto svoji poučku
nekompromisně.

Mohl by se jevit jako jakýsi \textit{,,hyperpelagianista``}. Vždyť neguje prvek
Božího výběru a Boží svobody v~lidské spáse. Predestinaci unifikuje na
univerzální predestinaci ke spasení a jediné, na čem záleží, je lidský prvek.
Zde je však nutné vzít v~potaz, že to jediné, co člověk může udělat, aby se
disponoval spáse, je podle Makoně odstoupení od sebe. Mnohokrát zdůrazňuje, že
vlastní silou se člověk nikdy nemůže dostat dál než k~extázi, kterou zavrhuje
jako překonaný fenomén.\rref{78-05A-Kaly}{07:12} I pro Makoně vše záleží na Bohu, který nám však dal
svobodu bránit se spáse.

Mohli bychom ho obvinit z~ezoterismu. Vždyť hovoří o rozmlouvání se záhrobním
světem,\rref{kotouc-D01-d}{08:10} o vílách\rref{80-07B}{24:46} a o
magii.\rref{80-12}{00:00} Nikdy však obcování s~\textit{,,vyššími sférami``}
nemá za cíl a nevede k~nim. Hovoří o nich buď pro odlehčení nebo častěji pro
dokreslení kontextu, vždy s~jasným cílem pomoci člověku co nejvíce na cestě do
Království Božího.

Jeho práce s~Biblí je velmi svérázná a podle mě ji nelze napodobit či z~ní
učinit metodu, ačkoliv se jeví, že Makoň po svých posluchačích právě toto
chtěl.\footnote{\rid{kotouc-77-brezen01-d}{01:25:56}, \rid{82-07}{01:01:00}}
Vychází někdy z~pochybných či zcela mylných stanovisek. Například že Máří
Magdaléna je totožná s~prostitutkou odsouzenou k~ukamenování (J 8, 3-11),\rref{83-22B}{29:46} nebo že jsme
všichni potomci Kainovi,\rref{83-27}{31:09} nebo že \textit{,,Bůh tak svět miloval, že dal svého jediného
Syna, aby žádný, kdo v~něho věří, nezahynul, ale měl věčný život,``} (J 3, 16) jsou slova
svatého Pavla.\rref{80-01B}{12:24} Několikrát tvrdí, že dokáže vyložit jakékoliv
podobenství,\rref{86-25A}{27:20} ovšem
prakticky se omezuje jen na hrstku z~nich. Radí vidět Bibli jako dokonalý návod
na cestu do Království Božího rozvedený do všech detailů, ovšem při svých
exegezích vrhá příběhy do jiného světla, než v~jakém jsou v~samotné
Bibli.\rref{85-06B}{24:43}

Zajisté to má souvislost s~tím, že dostal filozofické a ekonomické vzdělání a
nikoliv teologické a biblistické. Makoň křesťanství zná z~výchovy v~tradičním
římskokatolickém prostředí a detaily dostudovával později z~dostupných zdrojů na
vlastní pěst. Znalosti má rozsáhlé, ale zcela pochopitelně mu schází moderní
poznatky v~biblistice.

Pro mne osobně toto bylo asi největším rozčarováním, neboť jsem dlouho Makoňova
slova měl za ryzí pravdu a cestu, kterou lze doslova následovat a splnit tak
výzvu, kterou před nás předkládá: Vejít v~rámci tohoto života do vědomého
spojení s~Věčností, do Království Božího. Stále tuto jeho výzvu chci přijmout,
ovšem už si nemyslím, že by stačilo se řídit Makoňovou naukou coby jediným
vodítkem. Makoňovu práci s~Biblí převzít nedokážu a ani nechci. Stěžejní přístup
ale ano, a nemohu se tak vzdát svého vlastního přínosu, svojí vlastní cesty,
která se nutně od té Makoňovy musí lišit. Budiž proto Bohu vzdána čest a dík za
to, že Makoňova nauka má i své nedostatky. Ty jsou však vůči hodnotě, kterou
Makoňovo dílo má, naprosto nicotné.

Karel Makoň nám odkázal něco, co palčivě chybí: Životní program s~nejvyšším
možným cílem, strhávající celé naše mnohovrstevnaté a roztěkané osobnosti, a využívající
všechno, co svět nabízí. Nezavrhuje nic. Tradici, vědu, technologii, utrpení,
požitky, zkrátka všechno je pro něho prostředkem na cestě k Bohu, a žádného
prostředku se na této cestě nemáme štítit.

Přidávám se k~autorům Mozaiky mystického života Karla Makoně v~hluboké vděčnosti
za to, že jsem se s~jeho odkazem mohl setkat. Kéž tento skvost prospěje mnohým
dalším.


%%% Seznam použité literatury
%%% Seznam použité literatury (bibliografie)
%%%
%%% Pro vytváření bibliografie používáme bibTeX. Ten zpracovává
%%% citace v textu (např. makro \cite{...}) a vyhledává k nim literaturu
%%% v souboru literatura.bib.
%%%
%%% Příkaz \bibliographystyle určuje, jakým stylem budou citovány odkazy
%%% v textu. V závorce je název zvoleného souboru .bst. Styly plainnat
%%% a unsrt jsou standardní součástí latexových distribucí. Styl czplainnat
%%% je dodáván s touto šablonou a bibTeX ho hledá v aktuálním adresáři.

\bibliographystyle{czplainnat}    %% Autor (rok) s českými spojkami
% \bibliographystyle{plainnat}    %% Autor (rok) s anglickými spojkami
% \bibliographystyle{unsrt}       %% [číslo]

\renewcommand{\bibname}{Seznam použité literatury}

%%% Vytvoření seznamu literatury. Pozor, pokud jste necitovali ani jednu
%%% položku, seznam se automaticky vynechá.

\bibliography{citace}

%%% Kdybyste chtěli bibliografii vytvářet ručně (bez bibTeXu), lze to udělat
%%% následovně. V takovém případě se řiďte normou ISO 690 a zvyklostmi v oboru.

% \begin{thebibliography}{99}
%
% \bibitem{lamport94}
%   {\sc Lamport,} Leslie.
%   \emph{\LaTeX: A Document Preparation System}.
%   2. vydání.
%   Massachusetts: Addison Wesley, 1994.
%   ISBN 0-201-52983-1.
%
% \end{thebibliography}


%%% Obrázky v disertační práci
%%% (pokud jich je malé množství, obvykle není třeba seznam uvádět)
%\listoffigures

%%% Tabulky v disertační práci (opět nemusí být nutné uvádět)
%%% U matematických prací může být lepší přemístit seznam tabulek na začátek práce.
%\listoftables

%%% Použité zkratky v disertační práci (opět nemusí být nutné uvádět)
%%% U matematických prací může být lepší přemístit seznam zkratek na začátek práce.
%\chapwithtoc{Seznam použitých zkratek}

%%% Součástí doktorských prací musí být seznam vlastních publikací
%\input{publikace.tex}

%%% Přílohy k disertační práci, existují-li. Každá příloha musí být alespoň jednou
%%% odkazována z vlastního textu práce. Přílohy se číslují.
%%%
%%% Do tištěné verze se spíše hodí přílohy, které lze číst a prohlížet (dodatečné
%%% tabulky a grafy, různé textové doplňky, ukázky výstupů z počítačových programů,
%%% apod.). Do elektronické verze se hodí přílohy, které budou spíše používány
%%% v elektronické podobě než čteny (zdrojové kódy programů, datové soubory,
%%% interaktivní grafy apod.). Elektronické přílohy se nahrávají do SISu a lze
%%% je také do práce vložit na CD/DVD. Povolené formáty souborů specifikuje
%%% opatření rektora č. 72/2017.
%\appendix
%\chapter{Přílohy}

%\section{První příloha}

\openright
\end{document}
